\begingroup
\renewcommand{\figura}[1]{\begin{center}\begin{adjustbox}{max width=\linewidth,max totalheight=4.05cm}#1\end{adjustbox}\end{center}}
\section{Fabricação}
\begin{frame}{Do inversor ao silício: o que precisa ser fabricado?}
\only<1|handout:0>{\figura{\corteCMOS{1}}\etapa{1/3}{Definir regiões do tipo P e N e isolar os dispositivos.}}
\only<2|handout:0>{\figura{\corteCMOS{2}}\etapa{2/3}{Formar o dielétrico e o Gate; criar Source e Drain.}}
\only<3|handout:1>{\figura{\corteCMOS{3}}\etapa{3/3}{Contatos e metais conectam os transistores para realizar o inversor.}}
\fonte{Corte conceitual de CMOS planar bulk, sem escala. Etapas agrupadas para explicar a estrutura; não é receita de processo. Óxidos e conexões de corpo simplificados.}
\end{frame}

\begin{frame}{Fabricar é repetir operações em muitas camadas}
\figura{\cicloFabricacao}
\chamada{A litografia define onde uma operação deve acontecer.}
\fonte{Síntese didática: \href{https://www.asml.com/en/technology/all-about-microchips/how-microchips-are-made}{ASML — How microchips are made}. Deposição, remoção e dopagem se combinam em sequências específicas de cada processo; não ocorrem todas em todo ciclo.}
\end{frame}

\begin{frame}{Litografia: desenhar padrões com luz}
\only<1|handout:0>{\figura{\litografiaEtapa{1}}\etapa{1/3}{Exposição: a luz altera a solubilidade do resiste.}}
\only<2|handout:0>{\figura{\litografiaEtapa{2}}\etapa{2/3}{Revelação: o padrão aparece no resiste.}}
\only<3|handout:1>{\figura{\litografiaEtapa{3}}\etapa{3/3}{Transferência: gravar a camada; depois remover o resiste.}}
\fonte{\href{https://www.asml.com/en/technology/lithography-principles}{ASML — Lithography principles}. Resiste positivo; luz no wafer após projeção óptica. Máscara e raios simplificados, sem escala. O ciclo se repete para diferentes camadas.}
\end{frame}

\begin{frame}{A fronteira da litografia: DUV, EUV e High-NA}
\figura{\begin{tikzpicture}[x=1cm,y=1cm,font=\small]
\foreach \x/\titulo/\onda/\sub in {0/DUV/193/Lentes · imersão,3.8/EUV/{13{,}5}/Espelhos · vácuo,7.6/{EUV High-NA}/{13{,}5}/Óptica de maior abertura}{%
\node[blue_theme,font=\small\bfseries] at (\x+1.2,4.15){\titulo};
\node at (\x+1.2,3.5){$\lambda=\onda$ nm};
\node[align=center,font=\scriptsize,text width=3.3cm] at (\x+1.2,2.9){\sub};}
% Cones ilustram abertura numérica, não uma máquina real.
\fill[blue_theme!12] (4.3,2.2)--(5,.7)--(5.7,2.2)--cycle;
\fill[blue_theme!22] (7.6,2.2)--(8.8,.7)--(10,2.2)--cycle;
\draw[blue_theme,thick] (4.3,2.2)--(5,.7)--(5.7,2.2);\draw[blue_theme,thick] (7.6,2.2)--(8.8,.7)--(10,2.2);
\draw[blue_theme,thick] (4,.65)--(6,.65) (7.6,.65)--(10,.65);
\draw[blue_theme,thick] (.5,.65)--(1.9,.65);\draw[blue_theme,thick] (.3,1.7) ellipse (.9 and .15);\draw[seta] (1.2,2.2)--(1.2,.8);
\node[align=center,font=\scriptsize] at (1.2,-.05){DUV segue em uso\\em muitas camadas};
\node[align=center,font=\scriptsize] at (5,-.05){NA = 0,33\\resolução: 13 nm};
\node[align=center,font=\scriptsize] at (8.8,-.05){NA = 0,55\\resolução: 8 nm};
\path[use as bounding box] (-.4,-.7) rectangle (10.6,4.6);
\end{tikzpicture}}
\chamada{Luz, resolução e nome do nó: três grandezas distintas.}
\fonte{ASML: \href{https://www.asml.com/en/technology/lithography-principles/light-and-lasers}{luz}; \href{https://www.asml.com/en/company/stories/2024/5-things-high-na-euv}{High-NA EUV}. Resolução nominal por exposição nas plataformas NXE/EXE; não é tamanho de transistor. NA: abertura numérica.}
\end{frame}

\endgroup
